\cotefile{34}
\scanpage{50}
\mspage{48}
\margin{(\unsure{Deligne})}
% lines 1-7
\inserted{Proposition} \deleted{\uline{Corollaire}} 9.6. Supposons résolution des singularités sur $K^{\infty}$ en ttes dimensions, et que pour tt $X$ \deleted{\unsure{qui}} \uline{projectif et lisse} sur $K$, \deleted{\unsure{les}} l'action de $I$ sur $\unsure{H^i_q|H^i_c}(X_{\overline K},\mathbb{Q}_\ell)$ soit quasi-unipotente, et \unsure{kif-kif} après extension finie radicielle de $K$. Alors \deleted{\unsure{l'action de}} pour tt $X_K$ de t.f. sur $K$, l'action de $I$ sur les $H^i(X_{\overline K},\mathbb{Q}_\ell)$ est quasi-unipotente.

% lines 8-9
Il suffit de \unsure{conjuguer|combiner} \deleted{\unsure{9.6}} \inserted{\unsure{ce|la} qui précède} et \unsure{9.4. c)}, en procédant par récurrence sur \unsure{$n$}.

% lines 10-14
\deleted{\uline{Corollaire} \unsure{9.6.}}

9.7. \uline{Cas des} $H^i(X_{\overline K},\mathbb{Q}_\ell)$, pour $i\leq 1$.

Le cas de $i=0$ est trivial et laissé au lecteur (en fait, $\pi$ y opère de façon \uuline{quasi-triviale}).

% lines 15-21
Pour le cas $i=1$, on \deleted{\unsure{se ramène}} peut, quitte à passer à une extension finie de $K$, supposer que les composantes irréductibles réduites de $X_K$ sont géom. intègres. Supposons $X$ réduit, soit $X'$ sa normalisée, $X''=X'\times_X X'$, \inserted{$X'''=X'\times_X X'\times_X X'$}, \unsure{etc.} Alors, on \deleted{\unsure{a}} a une suite spectrale de descente finie (SGA 4 \uline{VIII}) :
% lines 22-24
\[
\deleted{0\longrightarrow\operatorname{Ker}\bigl(H^0(X''_{\overline K})\overset{\partial}{\longrightarrow}H^0(X'''_{\overline K})\bigr)}
\]
\[
\begin{aligned}
0\longrightarrow H^1\bigl(\alpha\longmapsto H^0((\unsure{X'|X^\bullet}/X)^{\alpha+1}_{\overline K},\mathbb{Q}_\ell)\bigr)&\longrightarrow H^1(X_{\overline K},\mathbb{Q}_\ell)\deleted{\unsure{\longrightarrow H^1}}\\
&\longrightarrow\operatorname{Ker}\bigl(H^1(X'_{\overline K},\mathbb{Q}_\ell)\overset{\partial}{\underset{\unsure{\neq|\delta}}{\longrightarrow}}H^1(X''_{\overline K},\mathbb{Q}_\ell)\bigr)
\end{aligned}
\]
\margin{\unsure{quotient d'un sous-espace} de $H^0(X''_{\overline K},\mathbb{Q}_\ell)$}
% lines 25-26
Cela nous ramène donc à voir que l'action de $I$ sur \deleted{\unsure{$X'$}} $H^1(X'_{\overline K},\mathbb{Q}_\ell)$. Donc on est ramené au
% archive sheet number: 49
